\begin{definition}[Overlapping two-site supports]\label{def:overlapping_two_site_supports}
    \lean{MPSTensor.axPairCfg}
    \lean{MPSTensor.xbPairCfg}
    \lean{MPSTensor.replaceAXCfg}
    \lean{MPSTensor.replaceXBCfg}
    \lean{finThreeArrowEquiv}
    \lean{MPSTensor.leftPairLift}
    \lean{MPSTensor.rightPairLift}
    \lean{MPSTensor.leftPairLift_sub}
    \lean{MPSTensor.rightPairLift_sub}
    \lean{MPSTensor.leftPairLift_one_sub}
    \lean{MPSTensor.rightPairLift_one_sub}
    \lean{MPSTensor.HasOverlappingTwoSiteCommutation}
    \lean{MPSTensor.HasAppendixD2ParentCommutingHamiltonian}
    \lean{MPSTensor.HasAppendixD2ParentCommutingHamiltonian.to_overlapping}
    \lean{MPSTensor.HasOverlappingTwoSiteCommutation.left_complement_idempotent}
    \lean{MPSTensor.HasOverlappingTwoSiteCommutation.right_complement_idempotent}
    \lean{MPSTensor.HasOverlappingTwoSiteCommutation.commute_complement_lifts}
    \lean{MPSTensor.HasOverlappingTwoSiteCommutation.complement}
    \lean{MPSTensor.adjacent_twoSite_cyclicWindowsOverlap}
    \leanok
    \uses{def:nsite_space_parent, def:cyclic_window_support,
        def:cyclic_windows_overlap}
    On a three-site space
    $\mc H_A\otimes\mc H_X\otimes\mc H_B$, let $Q_{AX}$ and $Q_{XB}$ be
    two-site projectors. Their local actions on the three-site space are the
    operators
    \begin{align}
        (Q_{AX}\psi)(a,x,b)
        &=
        Q_{AX}(\alpha,\xi\mapsto\psi(\alpha,\xi,b))(a,x),
        \notag\\
        (Q_{XB}\psi)(a,x,b)
        &=
        Q_{XB}(\xi,\beta\mapsto\psi(a,\xi,\beta))(x,b).
        \notag
    \end{align}
    The algebraic part of the overlapping nearest-neighbor commutation
    condition is
    \begin{align}
        Q_{AX}^2&=Q_{AX},
        \notag\\
        Q_{XB}^2&=Q_{XB},
        \notag\\
        Q_{AX}Q_{XB}&=Q_{XB}Q_{AX},
        \notag
    \end{align}
    after these two local actions are viewed on the common three-site space.
    This is the local-support part of the parent commuting Hamiltonian condition
    in \cite[Definition~D.2]{Cirac2016MPDO_arXiv}. The cyclic windows beginning
    at $i$ and $i+1$ realize the two supports $AX$ and $XB$, and hence overlap at
    the middle site.

    The source local condition also fixes a three-site subspace $K_{AXB}$ and
    imposes the kernel-intersection equation
    $K_{AXB}=\ker((Q_{AX})_{AX})\cap\ker((Q_{XB})_{XB})$.
    Forgetting this intersection equation gives the algebraic overlapping
    two-site commutation condition displayed above.

    The two lift maps also distribute over complementary projectors:
    \begin{align}
        (1-Q_{AX})_{AX}&=1-(Q_{AX})_{AX},
        \notag\\
        (1-Q_{XB})_{XB}&=1-(Q_{XB})_{XB}.
        \notag
    \end{align}
    Consequently, if the projectors $Q_{AX}$ and $Q_{XB}$ commute on the
    three-site space, then their complementary projectors
    $P_{AX}=1-Q_{AX}$ and $P_{XB}=1-Q_{XB}$ commute after the same lifts.
    This is the elementary algebraic passage used when one moves between the
    $Q$-projectors of
    \cite[Appendix~D.2]{Cirac2016MPDO_arXiv} and the complementary ground-space
    projectors $P$. Since the original two projectors are idempotent, the
    complementary projectors are idempotent as well; hence the pair
    $P_{AX},P_{XB}$ satisfies the same idempotent-and-commutation condition on
    the three-site space.
\end{definition}

\begin{theorem}[Matrices of the two adjacent factor actions]
    \label{thm:pair_lift_kronecker_matrices}
    \lean{MPSTensor.leftPairLift_toMatrix_reindex}
    \lean{MPSTensor.rightPairLift_toMatrix_reindex}
    \leanok
    \uses{def:overlapping_two_site_supports}
    Let $\mc H=\C^d$, and let $Q$ be an endomorphism of
    $\mc H\otimes\mc H$. Under the canonical identifications of two- and
    three-site configurations with pairs and right-associated triples, the
    matrices of its two adjacent actions are
    \begin{align}
        [Q_{AX}]&=[Q]\otimes I_d,
        \notag\\
        [Q_{XB}]&=I_d\otimes[Q].
        \notag
    \end{align}
    Equivalently, for all $a,x,b,a',x',b'\in\Fin d$,
    \begin{align}
        [Q_{AX}]_{(a,x,b),(a',x',b')}
        &=[Q]_{(a,x),(a',x')}\delta_{b,b'},
        \notag\\
        [Q_{XB}]_{(a,x,b),(a',x',b')}
        &=\delta_{a,a'}[Q]_{(x,b),(x',b')}.
        \notag
    \end{align}
    This is the homogeneous three-single-site specialization of the factor
    actions on $\mc H_A\otimes\mc H_X\otimes\mc H_B$ in
    \cite[Definition~D.2]{Cirac2016MPDO_arXiv}; here
    $\mc H_A=\mc H_X=\mc H_B=\mc H$. No assertion is made for unequal factor
    spaces.
\end{theorem}

\begin{proof}
    \leanok
    On the first pair, the final coordinate is unchanged, which contributes
    $\delta_{b,b'}$ to the matrix coefficient. On the final pair, the first
    coordinate is unchanged, which contributes $\delta_{a,a'}$.
\end{proof}

\begin{theorem}[Three-site translated local terms]
    \label{thm:local_term_two_three_ax_xb_lifts}
    \lean{MPSTensor.localTerm_two_three_zero_eq_leftPairLift_parentInteraction}
    \lean{MPSTensor.localTerm_two_three_one_eq_rightPairLift_parentInteraction}
    \leanok
    \uses{def:overlapping_two_site_supports, def:local_term_parent}
    Put $q(A)=1-P_{G_2(A)}$ for the canonical two-site parent interaction. On
    the three-site space
    $\mc H_A\otimes\mc H_X\otimes\mc H_B$, with sites $0,1,2$ labelled
    $A,X,B$, the two adjacent translated length-two parent interactions are
    the two local actions
    \begin{align}
        h^{(3)}_0(A,2)&=q(A)_{AX},
        \notag\\
        h^{(3)}_1(A,2)&=q(A)_{XB}.
        \notag
    \end{align}
    This identifies the two translated parent interactions with the two
    supports appearing in \cite[Definition~D.2]{Cirac2016MPDO_arXiv}; it does
    not assert that they commute.
\end{theorem}

\begin{proof}\leanok
    Let $\sigma=(a,x,b)$ be a three-site configuration. For the first
    translated term, the length-two window and the corresponding replacement
    satisfy
    \begin{align}
        \sigma_{\{0,1\}}&=(a,x)=\sigma_{AX},
        \notag\\
        \sigma^{\{0,1\}\leftarrow(\alpha,\xi)}
        &=(\alpha,\xi,b)
        =\sigma^{AX\leftarrow(\alpha,\xi)}.
        \notag
    \end{align}
    For the second translated term,
    \begin{align}
        \sigma_{\{1,2\}}&=(x,b)=\sigma_{XB},
        \notag\\
        \sigma^{\{1,2\}\leftarrow(\xi,\beta)}
        &=(a,\xi,\beta)
        =\sigma^{XB\leftarrow(\xi,\beta)}.
        \notag
    \end{align}
    Substituting these two pairs of identities into the definition of the
    translated parent term gives the two displayed equations.
\end{proof}

\begin{theorem}[Three-site commutation from overlapping supports]
    \label{thm:local_term_two_three_commute_from_ax_xb}
    \lean{MPSTensor.localTerm_two_three_zero_one_commute_of_overlapping_two_site_commutation}
    \leanok
    \uses{def:overlapping_two_site_supports,
        thm:local_term_two_three_ax_xb_lifts}
    Suppose the canonical two-site parent interaction $q(A)=1-P_{G_2(A)}$
    satisfies the overlapping $AX/XB$ commutation condition of
    Definition~\ref{def:overlapping_two_site_supports}. Then, on the three-site
    chain,
    $h^{(3)}_0(A,2)h^{(3)}_1(A,2) = h^{(3)}_1(A,2)h^{(3)}_0(A,2)$.
    This only identifies the first two translated local terms with the lifted
    $AX$ and $XB$ actions; the construction of the
    \cite[Appendix~D.2]{Cirac2016MPDO_arXiv} projectors associated with the
    \cite[Appendix~B]{Cirac2016MPDO_arXiv} basic-vector form is a separate step.
\end{theorem}

\begin{proof}\leanok
    By Theorem~\ref{thm:local_term_two_three_ax_xb_lifts}, the two translated
    local terms are $q(A)_{AX}$ and $q(A)_{XB}$. The assumed overlapping
    support condition says exactly that these two lifted endomorphisms commute.
\end{proof}

\begin{theorem}[Three-site commutation from two-site representatives]
    \label{thm:local_term_two_three_commute_from_two_site_q}
    \lean{MPSTensor.localTerm_two_three_zero_one_commute_of_overlapping_two_site_projectors}
    \leanok
    \uses{def:overlapping_two_site_supports,
        thm:local_term_two_three_ax_xb_lifts}
    Let $Q_{AX}$ and $Q_{XB}$ be two-site projectors whose lifted $AX$ and
    $XB$ actions commute. Suppose that $q(A)=Q_{AX}$ and $q(A)=Q_{XB}$.
    Then the first two translated length-two local terms on the three-site
    chain commute.
\end{theorem}

\begin{proof}\leanok
    The lift identities identify the translated terms with the lifted actions
    of the two projectors:
    \begin{align}
        h^{(3)}_0(A,2)&=(Q_{AX})_{AX},
        \notag\\
        h^{(3)}_1(A,2)&=(Q_{XB})_{XB}.
        \notag
    \end{align}
    The overlapping-support hypothesis is the equation
    $(Q_{AX})_{AX}(Q_{XB})_{XB} = (Q_{XB})_{XB}(Q_{AX})_{AX}$.
    Substituting the two lifted identities gives the desired commutation of
    $h^{(3)}_0(A,2)$ and $h^{(3)}_1(A,2)$.
\end{proof}

\begin{theorem}[Three-site commutation from source projectors]
    \label{thm:local_term_two_three_commute_from_source_q}
    \lean{MPSTensor.localTerm_two_three_zero_one_commute_of_appendixD2}
    \leanok
    \uses{def:overlapping_two_site_supports,
        thm:local_term_two_three_ax_xb_lifts}
    Suppose the projectors $Q_{AX}$ and $Q_{XB}$ satisfy the
    \cite[Definition~D.2]{Cirac2016MPDO_arXiv} local parent-commuting
    condition on $A X B$, and suppose that the canonical two-site parent
    interaction is identified, as a two-site operator, with the two source
    projectors: $q(A)=Q_{AX}$ and $q(A)=Q_{XB}$. After lifting these two
    equalities to the $AX$ and $XB$ faces, the first two translated length-two
    local terms on the three-site chain commute. Thus,
    $h^{(3)}_0(A,2)h^{(3)}_1(A,2) = h^{(3)}_1(A,2)h^{(3)}_0(A,2)$.
\end{theorem}

\begin{proof}\leanok
    The lift identities of Theorem~\ref{thm:local_term_two_three_ax_xb_lifts}
    identify the translated terms with $(Q_{AX})_{AX}$ and $(Q_{XB})_{XB}$.
    After forgetting the kernel-intersection equation, the
    \cite[Definition~D.2]{Cirac2016MPDO_arXiv} condition says exactly that
    these two lifted projectors commute.
\end{proof}

\begin{theorem}[Three-site commutation from complementary Appendix projectors]
    \label{thm:local_term_two_three_commute_from_complement_ax_xb}
    \lean{MPSTensor.localTerm_two_three_zero_one_commute_of_overlapping_two_site_complement}
    \lean{MPSTensor.localTerm_two_three_zero_one_commute_of_appendixD2_complement}
    \leanok
    \uses{def:overlapping_two_site_supports,
        thm:local_term_two_three_ax_xb_lifts}
    Suppose the Appendix~D.2 projectors $Q_{AX}$ and $Q_{XB}$ satisfy the
    local parent-commuting condition on $A X B$, and suppose that the canonical
    two-site parent interaction is their complement on both faces:
    $q(A)=1-Q_{AX}=1-Q_{XB}$. Then the first two translated length-two local
    terms on the three-site chain commute,
    $h^{(3)}_0(A,2)h^{(3)}_1(A,2) = h^{(3)}_1(A,2)h^{(3)}_0(A,2)$.
\end{theorem}

\begin{proof}\leanok
    The lift identities of Theorem~\ref{thm:local_term_two_three_ax_xb_lifts}
    identify the translated terms with $(1-Q_{AX})_{AX}$ and
    $(1-Q_{XB})_{XB}$. The Appendix~D.2 condition implies the overlapping
    two-site commutation condition after forgetting the kernel-intersection
    equation. Definition~\ref{def:overlapping_two_site_supports} then gives
    commutation of the complementary lifted projectors, hence the displayed
    commutation relation.
\end{proof}

\begin{definition}[Commuting local terms for a parent Hamiltonian]\label{def:is_commuting_parent_ham}
    \lean{MPSTensor.IsCommutingParentHam}
    \leanok
    \uses{def:local_term_parent}
    The parent Hamiltonian with block length~$L$ on $N$ sites has
    \emph{commuting} local terms if $h_i h_j=h_j h_i$ for all
    $i,j\in\Fin N$. This records the translated commutation equation
    $[\tau_j(P_L),P_L]=0$ from
    \cite[Definition~3.9]{Cirac2016MPDO_arXiv}; the same source definition
    also requires the ground space to be spanned, for $N>L$, by the periodic
    MPS vectors of the BNT components.
\end{definition}

\begin{definition}[Nearest-neighbor commuting parent Hamiltonian]\label{def:is_nncph}
    \lean{MPSTensor.IsNNCPH}
    \leanok
    \uses{def:is_commuting_parent_ham}
    The nearest-neighbor commuting condition is the length-$2$ specialization
    of the preceding commutation equation. For every $i,j\in\Fin N$, one has
    $h_i h_j=h_j h_i$. The source definition
    \cite[Definition~3.9]{Cirac2016MPDO_arXiv} also requires the corresponding
    parent Hamiltonian to have $L=2$ and to have the prescribed ground-space
    spanning property for $N>2$.
\end{definition}

\begin{lemma}[Overlap reduction for commuting local terms]
    \label{lem:commuting_parent_ham_overlap_reduction}
    \lean{MPSTensor.isCommutingParentHam_of_cyclicWindowsOverlap_commute}
    \lean{MPSTensor.isNNCPH_of_twoSite_cyclicWindowsOverlap_commute}
    \leanok
    \uses{def:is_commuting_parent_ham, def:is_nncph,
        def:cyclic_windows_overlap, lem:transported_es_nonoverlap_positive}
    Let $L\le N$. Suppose that $h_i h_j=h_j h_i$ whenever the length-$L$
    cyclic supports of the two local terms meet. Then the parent Hamiltonian
    has commuting local terms: $h_i h_j=h_j h_i$ for all $i,j$. In the case
    $L=2$, this gives the nearest-neighbor commuting condition.
\end{lemma}

\begin{proof}\leanok
    \uses{lem:transported_es_nonoverlap_positive}
    If the two cyclic supports meet, the hypothesis gives the commutation
    relation. If they are disjoint, the corresponding local terms commute by
    locality.
\end{proof}

\begin{lemma}[Length-two cyclic-overlap alternatives]
    \label{lem:two_site_cyclic_window_overlap_cases}
    \lean{MPSTensor.cyclicWindowsOverlap_twoSite_cases}
    \leanok
    \uses{def:cyclic_windows_overlap}
    For $i,j\in\Z/N\Z$, if the two length-two cyclic supports starting at $i$
    and $j$ overlap, then
    \begin{align}
        j&=i,
        \notag\\
        \text{or}\qquad j&\equiv i+1\pmod N,
        \notag\\
        \text{or}\qquad i&\equiv j+1\pmod N.
        \label{eq:parent_two_site_overlap_cases}
    \end{align}
\end{lemma}

\begin{proof}
    \leanok
    Expanding each length-two cyclic support gives shifts
    $r_i,r_j\in\{0,1\}$ and an equality of sites
    $i+r_i\equiv j+r_j\pmod N$. The four possible pairs $(r_i,r_j)$ give
    exactly the three alternatives in
    (\ref{eq:parent_two_site_overlap_cases}).
\end{proof}

\begin{lemma}[Adjacent two-site translates suffice for nearest-neighbor commutation]
    \label{lem:nncph_from_adjacent_two_site_commutation}
    \lean{MPSTensor.isNNCPH_of_adjacent_twoSite_commute}
    \lean{MPSTensor.HasProductPairLocalProjectors.of_adjacent_twoSite_commute}
    \leanok
    \uses{lem:two_site_cyclic_window_overlap_cases,
        lem:commuting_parent_ham_overlap_reduction, lem:local_term_idempotent}
    Let $N\ge2$. Suppose that, for every $i\in\Z/N\Z$, the translated
    length-two local term commutes with its cyclic right neighbor:
    \begin{align}
        h_i(A,2)h_{i+1}(A,2)
        &=
        h_{i+1}(A,2)h_i(A,2).
        \label{eq:parent_adjacent_two_site_commutation}
    \end{align}
    Then all translated length-two local terms commute, and the family
    $p_i=h_i(A,2)$ satisfies the finite-chain local-projector hypotheses used
    in the conditional Appendix~B extraction.
\end{lemma}

\begin{proof}
    \leanok
    By Lemma~\ref{lem:two_site_cyclic_window_overlap_cases}, overlapping
    length-two cyclic supports force the three alternatives in
    (\ref{eq:parent_two_site_overlap_cases}). The equal-start case gives
    $h_i(A,2)=h_j(A,2)$. The two nearest-neighbor cases are
    (\ref{eq:parent_adjacent_two_site_commutation}) and its transpose. The
    overlap-reduction lemma supplies the remaining disjoint pairs by locality.
\end{proof}

\begin{lemma}[Local-projector hypotheses from overlapping two-site commutation]
    \label{lem:product_pair_local_projectors_from_overlap}
    \lean{MPSTensor.HasProductPairLocalProjectors.of_twoSite_cyclicWindowsOverlap_commute}
    \leanok
    \uses{lem:commuting_parent_ham_overlap_reduction, lem:local_term_idempotent}
    Let $N\ge2$. Suppose that the translated length-two local terms commute for
    every pair of overlapping cyclic windows:
    $h_i(A,2)h_j(A,2) = h_j(A,2)h_i(A,2)$.
    Then the family $p_i=h_i(A,2)$ gives the finite-chain local-projector
    hypotheses used in the conditional Appendix~B extraction.
\end{lemma}

\begin{proof}\leanok
    Lemma~\ref{lem:commuting_parent_ham_overlap_reduction} turns the
    overlapping-pair commutation equations into commutation for all pairs. The
    idempotency equation is Lemma~\ref{lem:local_term_idempotent}.
\end{proof}

\begin{definition}[Commutation and zero-energy equations for nearest-neighbor local terms]%
    \label{def:is_nncph_ground_state}
    \lean{MPSTensor.IsNNCPHGroundState}
    \leanok
    \uses{def:is_nncph, def:is_frustration_free}
    This condition records the commutation equation and the zero-energy equation
    for the nearest-neighbor local terms:
    \begin{align}
        h_i h_j&=h_j h_i,
        \notag\\
        h_i V^{(N)}(A)&=0.
        \notag
    \end{align}
    This is the zero-energy part of the condition appearing in
    \cite[Theorem~3.10(iii)]{Cirac2016MPDO_arXiv}: the MPS vector has zero
    energy for the translated two-site local terms, and those terms commute.
    The missing source assertion is that the whole parent-Hamiltonian ground
    space is spanned by the corresponding periodic MPS vectors from the BNT
    components, as recorded in
    Definition~\ref{def:parent_hamiltonian_ground_space_spanning}.
    % Scope restriction recorded in
    % docs/paper-gaps/cpsv16_nncph_ground_state_scope.tex.
\end{definition}

\begin{definition}[Parent-Hamiltonian ground-space spanning condition]%
    \label{def:parent_hamiltonian_ground_space_spanning}
    \lean{MPSTensor.HasParentHamiltonianGroundSpaceSpanning}
    \leanok
    \uses{def:parent_hamiltonian, def:bnt_span}
    Let $B$ be a canonical-form tensor with BNT components
    $A_1,\ldots,A_g$. The source parent-Hamiltonian spanning condition is that,
    for every $N>L$,
    $\ker H_L^{(N)}(B)=\spn\{V^{(N)}(A_j):j=1,\ldots,g\}$.
    This is the ground-space clause in
    \cite[Definition~3.9]{Cirac2016MPDO_arXiv}. It is separate from the
    commutation equation $[\tau_j(P_L),P_L]=0$ and from the zero-energy
    equation for a single periodic MPS vector.
\end{definition}

\begin{definition}[Nearest-neighbor commuting parent-Hamiltonian ground-space condition]%
    \label{def:has_nncph_ground_space}
    \lean{MPSTensor.HasNNCPHGroundSpace}
    \leanok
    \uses{def:is_nncph_ground_state, def:parent_hamiltonian_ground_space_spanning}
    For a fixed chain length $N$, this condition is the conjunction of the
    nearest-neighbor commutation equation, the zero-energy equation for the MPS
    vector, and the ground-space spanning equation:
    \begin{align}
        h_i h_j&=h_j h_i,
        \notag\\
        h_iV^{(N)}(B)&=0,
        \notag\\
        \ker H_2^{(N)}(B)
        &=
        \spn\{V^{(N)}(A_j):j=1,\ldots,g\}.
        \notag
    \end{align}
    The source statement in
    \cite[Theorem~3.10(iii)]{Cirac2016MPDO_arXiv} quantifies this condition
    over all $N>2$.
\end{definition}

\begin{definition}[All-chain nearest-neighbor commuting parent-Hamiltonian ground-space condition]%
    \label{def:has_nncph_ground_spaces}
    \lean{MPSTensor.HasNNCPHGroundSpaces}
    \leanok
    \uses{def:has_nncph_ground_space}
    This is the source quantification in
    \cite[Theorem~3.10(iii)]{Cirac2016MPDO_arXiv}. For every $N>2$, the
    fixed-chain condition of Definition~\ref{def:has_nncph_ground_space}
    holds:
    \begin{align}
        h_i h_j&=h_j h_i,
        \notag\\
        h_iV^{(N)}(B)&=0,
        \notag\\
        \ker H_2^{(N)}(B)
        &=
        \spn\{V^{(N)}(A_j):j=1,\ldots,g\}.
        \notag
    \end{align}
    This condition records the ground-space spanning equation, but does not
    prove it for a concrete canonical-form tensor.
\end{definition}

\begin{remark}[Elementary consequences of the commuting definition]
    \label{rem:nncph_ground_spaces_decomposition}
    \lean{MPSTensor.IsNNCPH.isCommutingParentHam}
    \lean{MPSTensor.IsNNCPHGroundState.isNNCPH}
    \lean{MPSTensor.IsNNCPHGroundState.isFrustrationFree}
    \lean{MPSTensor.IsNNCPH.isNNCPHGroundState}
    \lean{MPSTensor.HasNNCPHGroundSpace.isNNCPHGroundState}
    \lean{MPSTensor.HasNNCPHGroundSpace.isNNCPH}
    \lean{MPSTensor.HasNNCPHGroundSpace.isFrustrationFree}
    \lean{MPSTensor.HasNNCPHGroundSpace.groundSpaceSpanning}
    \lean{MPSTensor.HasNNCPHGroundSpace.hasParentHamiltonianGroundSpaceSpanning}
    \lean{MPSTensor.HasNNCPHGroundSpaces.hasNNCPHGroundSpace}
    \lean{MPSTensor.HasNNCPHGroundSpaces.hasParentHamiltonianGroundSpaceSpanning}
    \lean{MPSTensor.hasNNCPHGroundSpaces_iff_forall_isNNCPHGroundState_and_groundSpaceSpanning}
    \lean{MPSTensor.hasNNCPHGroundSpaces_iff_forall_isNNCPH_and_groundSpaceSpanning}
    \lean{MPSTensor.IsCommutingParentHam.symm}
    \lean{MPSTensor.IsCommutingParentHam.ham_comm_localTerm}
    \leanok
    \uses{lem:parent_hamiltonian_ff, def:has_nncph_ground_space}
    The nearest-neighbor condition is the length-$2$ specialization of the
    general commuting condition. The ground-vector condition above contains
    this commuting equation together with the frustration-free equation for
    $V^{(N)}(A)$; it does not assert the source ground-space spanning property.
    For $N\ge2$, Lemma~\ref{lem:parent_hamiltonian_ff} with $L=2$ gives
    $h_iV^{(N)}(A)=0$ for every site $i$. The fixed-chain ground-space
    condition also contains
    $\ker H_2^{(N)}(B) = \spn\{V^{(N)}(A_j):j=1,\ldots,g\}$,
    and quantifying this condition over all $N>2$ gives the nearest-neighbor
    case of the spanning condition above.
    Conversely, the all-chain condition is exactly the conjunction of the
    all-chain ground-vector equations and the nearest-neighbor instance of the
    spanning condition above. This is precisely the condition named in
    Definition~\ref{def:has_nncph_ground_spaces}.
    Since the displayed zero-energy equation follows from the
    parent-Hamiltonian construction for $N>2$, independently of the
    commutation and spanning equations, this is equivalently
    \begin{align}
        &\forall N>2,\quad
        \left[\begin{array}{l}
        \forall\, i,j,\quad
        h_i^{(N)}h_j^{(N)}=h_j^{(N)}h_i^{(N)},\\
        \forall\, i,\quad
        h_i^{(N)}V^{(N)}(B)=0,\\
        \ker H_2^{(N)}(B)
            =
            \spn\{V^{(N)}(A_j):j=1,\ldots,g\}
        \end{array}\right]
        \notag\\
        &\qquad\Longleftrightarrow\qquad
        \left[\begin{array}{l}
        \forall N>2,\quad
        \forall\, i,j,\quad
        h_i^{(N)}h_j^{(N)}=h_j^{(N)}h_i^{(N)},\\
        \forall N>2,\quad
        \ker H_2^{(N)}(B)
            =
            \spn\{V^{(N)}(A_j):j=1,\ldots,g\}.
        \end{array}\right].
        \notag
    \end{align}
\end{remark}

\begin{theorem}[Chain ground-space vectors are frustration-free]
    \label{thm:chain_ground_space_implies_frustration_free}
    \lean{MPSTensor.isFrustrationFree_of_mem_chainGroundSpace}
    \leanok
    \uses{def:chain_ground_space, def:is_frustration_free}
    Assume $N>0$ and $L\le N$. Let $\psi$ be a vector in the periodic chain
    ground space $G_L^{(N)}(A)$. Then $\psi$ satisfies all local zero-energy
    equations: for every site $i$, one has $h_i(A,L)\psi=0$.
\end{theorem}

\begin{proof}
    \leanok
    Write $R_{i,\eta}\psi$ for the cyclic length-$L$ restriction of $\psi$ at
    site $i$, with the outside configuration fixed by $\eta$. Membership in
    $G_L^{(N)}(A)$ gives
    $R_{i,\eta}\psi\in G_L(A)=\ker h_L(A)$ for every $i$ and $\eta$. Hence
    $h_L(A)R_{i,\eta}\psi=0$. For a full configuration $\sigma$, taking
    $\eta=\sigma|_{\{i,\ldots,i+L-1\}^c}$ gives
    \begin{align}
        (h_i(A,L)\psi)(\sigma)
        &=
        (h_L(A)R_{i,\eta}\psi)
            (\sigma|_{\{i,\ldots,i+L-1\}}) =0.
        \notag
    \end{align}
    Thus $h_i(A,L)\psi=0$ for every $i$.
\end{proof}
